
% ==========================================
% PHAN D - CAU HINH BGP
% ==========================================

\section{Phần D -- Cấu hình BGP}

% ==========================================
% A. MUC TIEU
% ==========================================

\subsection{Mục tiêu}

Phần này triển khai giao thức Border Gateway
Protocol (BGP) để thực hiện định tuyến giữa
các Autonomous System (AS) trong mô hình
Internet thu nhỏ, đồng thời trao đổi thông tin
định tuyến giữa các router BGP thuộc cùng AS.

\textbf{Các yêu cầu kỹ thuật bao gồm:}

\begin{itemize}

    \item Thiết lập các phiên eBGP trên
    ba liên kết L4 (A2--B1), L5 (A3--C1)
    và L9 (B3--C1), sử dụng địa chỉ IP
    trực tiếp của các liên kết.

    \item Thiết lập iBGP giữa A2 và A3
    trong AS A, đồng thời cấu hình
    iBGP full mesh giữa B1, B2 và B3
    trong AS B.

    \item Sử dụng địa chỉ Loopback0
    để thiết lập các phiên iBGP.

    \item Xử lý thuộc tính NEXT\_HOP
    nhằm đảm bảo các router có thể
    truy cập địa chỉ next-hop của
    những tuyến BGP nhận được.

    \item Quảng bá chính xác bốn mạng
    LAN-A, LAN-B, LAN-C và LAN-D
    theo phân công của đề bài.

    \item Kiểm chứng trạng thái các
    phiên BGP và khả năng kết nối
    giữa sáu cặp mạng LAN.

\end{itemize}


% ==========================================
% B. TRIEN KHAI CAU HINH
% ==========================================

\subsection{Triển khai cấu hình}

\subsubsection{Nguyên tắc triển khai BGP}

BGP được triển khai theo hai hình thức:

\begin{itemize}

    \item \textbf{eBGP:}
    Thiết lập giữa các router thuộc
    những AS khác nhau, sử dụng
    địa chỉ IP của các liên kết trực tiếp.

    \item \textbf{iBGP:}
    Thiết lập giữa các router thuộc
    cùng một AS, sử dụng địa chỉ
    Loopback0 để tạo phiên BGP ổn định.

\end{itemize}

Các phiên eBGP cần thiết lập
được trình bày trong Bảng
\ref{tab:partD-ebgp}.

\begin{table}[H]
\centering
\caption{Các phiên eBGP giữa những AS}
\label{tab:partD-ebgp}
\small
\begin{tabularx}{\textwidth}{c c X X}
\toprule
\textbf{Link} &
\textbf{AS} &
\textbf{Router 1} &
\textbf{Router 2} \\
\midrule
L4 & A--B &
A2: 10.71.26.13 &
B1: 10.71.26.14 \\

L5 & A--C &
A3: 10.71.26.17 &
C1: 10.71.26.18 \\

L9 & B--C &
B3: 10.71.26.33 &
C1: 10.71.26.34 \\
\bottomrule
\end{tabularx}
\end{table}

Đối với iBGP, AS A có một phiên
A2--A3, trong khi AS B có ba phiên
B1--B2, B1--B3 và B2--B3.

Trước khi thiết lập iBGP,
các địa chỉ Loopback0 phải có
khả năng truy cập thông qua
giao thức định tuyến nội bộ.

Cụ thể, AS A sử dụng RIPv2
và AS B sử dụng OSPF đã được
triển khai trong các phần trước.


% ==========================================
% CAC LENH BGP CHINH
% ==========================================

\subsubsection{Giải thích các lệnh cấu hình BGP chính}

Các lệnh BGP được sử dụng trong
quá trình triển khai bao gồm:

\begin{itemize}

    \item \textbf{\texttt{router bgp <AS>}:}
    Khởi tạo tiến trình BGP với
    số Autonomous System tương ứng.

    \item \textbf{\texttt{neighbor <IP> remote-as <AS>}:}
    Khai báo địa chỉ IP của router
    láng giềng và số AS của router đó.
    Nếu hai router thuộc cùng AS,
    phiên được thiết lập là iBGP;
    nếu khác AS, phiên là eBGP.

    \item \textbf{\texttt{neighbor <IP> update-source Loopback0}:}
    Sử dụng địa chỉ Loopback0
    làm địa chỉ nguồn cho phiên BGP.
    Lệnh này được áp dụng cho
    các phiên iBGP trong đồ án.

    \item \textbf{\texttt{neighbor <IP> next-hop-self}:}
    Thay đổi thuộc tính NEXT\_HOP
    của các tuyến được quảng bá
    cho láng giềng BGP, sử dụng
    địa chỉ của chính router
    đang gửi cập nhật làm next-hop.

    Trong đồ án, lệnh này được sử dụng
    trên các phiên iBGP để xử lý
    những tuyến nhận từ AS bên ngoài.

    \item \textbf{\texttt{network <IP> mask <MASK>}:}
    Quảng bá một mạng vào BGP
    khi có tuyến khớp chính xác
    với địa chỉ mạng và subnet mask
    trong bảng định tuyến IP.

\end{itemize}

Ví dụ, câu lệnh:

\begin{lstlisting}
network 10.171.126.64 mask 255.255.255.192
\end{lstlisting}

được sử dụng để quảng bá mạng LAN-B
với prefix /26 vào BGP.

Lệnh \texttt{network} không tự tạo
một tuyến đường mới trong bảng
định tuyến IP. Mạng được quảng bá
phải có sẵn tuyến tương ứng
trong Routing Information Base (RIB).


% ==========================================
% CAU HINH AS A
% ==========================================

\subsubsection{Cấu hình BGP trong AS A}

AS A có hai router biên A2 và A3.
A2 thiết lập eBGP với B1 thuộc AS B,
trong khi A3 thiết lập eBGP với C1
thuộc AS C.

Đồng thời, A2 và A3 thiết lập
iBGP với nhau thông qua địa chỉ
Loopback0 để trao đổi thông tin
định tuyến nhận được từ hai ISP.

\textbf{Cấu hình trên router A2:}

\begin{lstlisting}
configure terminal

router bgp 64626

 neighbor 10.71.26.14 remote-as 64726

 neighbor 10.171.26.3 remote-as 64626
 neighbor 10.171.26.3 update-source Loopback0
 neighbor 10.171.26.3 next-hop-self

 network 10.171.126.0 mask 255.255.255.192

end
\end{lstlisting}

\textbf{Cấu hình trên router A3:}

\begin{lstlisting}
configure terminal

router bgp 64626

 neighbor 10.71.26.18 remote-as 64826

 neighbor 10.171.26.2 remote-as 64626
 neighbor 10.171.26.2 update-source Loopback0
 neighbor 10.171.26.2 next-hop-self

 network 10.171.126.0 mask 255.255.255.192

end
\end{lstlisting}

Hai router sử dụng địa chỉ
Loopback0 để thiết lập iBGP,
trong khi các phiên eBGP sử dụng
địa chỉ trên liên kết L4 và L5.

Lệnh \texttt{next-hop-self}
được áp dụng cho phiên iBGP
để các tuyến học từ ISP
có next-hop là router biên
tương ứng trong AS A.

Cả A2 và A3 đều quảng bá mạng
LAN-A \texttt{10.171.126.0/26}.
Tuyến đến LAN-A phải được học
thông qua RIPv2 và có mặt
trong bảng định tuyến IP
trước khi BGP quảng bá.


% ==========================================
% CAU HINH AS B
% ==========================================

\subsubsection{Cấu hình BGP trong AS B}

AS B bao gồm ba router B1, B2
và B3, được triển khai iBGP
theo mô hình full mesh.

Các phiên iBGP sử dụng Loopback0,
trong khi eBGP được thiết lập
trên B1 với A2 và trên B3 với C1.

\textbf{Cấu hình trên router B1:}

\begin{lstlisting}
configure terminal

router bgp 64726

 neighbor 10.71.26.13 remote-as 64626

 neighbor 10.171.26.5 remote-as 64726
 neighbor 10.171.26.5 update-source Loopback0
 neighbor 10.171.26.5 next-hop-self

 neighbor 10.171.26.6 remote-as 64726
 neighbor 10.171.26.6 update-source Loopback0
 neighbor 10.171.26.6 next-hop-self

end
\end{lstlisting}

Router B1 thiết lập eBGP với A2
thông qua L4 và iBGP với B2, B3
thông qua Loopback0.

\textbf{Cấu hình trên router B2:}

\begin{lstlisting}
configure terminal

router bgp 64726

 neighbor 10.171.26.4 remote-as 64726
 neighbor 10.171.26.4 update-source Loopback0
 neighbor 10.171.26.4 next-hop-self

 neighbor 10.171.26.6 remote-as 64726
 neighbor 10.171.26.6 update-source Loopback0
 neighbor 10.171.26.6 next-hop-self

 network 10.171.126.64 mask 255.255.255.192
 network 10.171.126.192 mask 255.255.255.192

end
\end{lstlisting}

B2 thiết lập hai phiên iBGP
với B1 và B3, đồng thời quảng bá
hai mạng LAN:

\begin{itemize}
    \item LAN-B: \texttt{10.171.126.64/26}.
    \item LAN-D: \texttt{10.171.126.192/26}.
\end{itemize}

LAN-B được kết nối trực tiếp
thông qua Loopback1 của B2.

Đối với LAN-D, B2 cần có
static route đến mạng này
thông qua D1, đã được triển khai
trong Phần A.

\textbf{Cấu hình trên router B3:}

\begin{lstlisting}
configure terminal

router bgp 64726

 neighbor 10.71.26.34 remote-as 64826

 neighbor 10.171.26.4 remote-as 64726
 neighbor 10.171.26.4 update-source Loopback0
 neighbor 10.171.26.4 next-hop-self

 neighbor 10.171.26.5 remote-as 64726
 neighbor 10.171.26.5 update-source Loopback0
 neighbor 10.171.26.5 next-hop-self

end
\end{lstlisting}

Router B3 thiết lập eBGP với C1
trên liên kết L9 và iBGP với
B1, B2 thông qua Loopback0.


% ==========================================
% CAU HINH AS C
% ==========================================

\subsubsection{Cấu hình BGP trong AS C}

AS C chỉ có một router C1,
đóng vai trò ISP kết nối với
AS A qua L5 và AS B qua L9.

Do AS C chỉ có một router,
không cần triển khai iBGP.

Cấu hình trên router C1:

\begin{lstlisting}
configure terminal

router bgp 64826

 neighbor 10.71.26.17 remote-as 64626
 neighbor 10.71.26.33 remote-as 64726

 network 10.171.126.128 mask 255.255.255.192
 network 10.171.126.192 mask 255.255.255.192

end
\end{lstlisting}

Router C1 quảng bá LAN-C
\texttt{10.171.126.128/26}
và LAN-D \texttt{10.171.126.192/26}.

LAN-C được kết nối trực tiếp
thông qua Loopback1 của C1,
trong khi LAN-D được truy cập
thông qua static route đến D1
đã cấu hình trong Phần A.


% ==========================================
% XU LY NEXT HOP
% ==========================================

\subsubsection{Xử lý thuộc tính NEXT\_HOP}

Khi một router biên nhận tuyến
từ eBGP và quảng bá tuyến đó
cho láng giềng iBGP, thuộc tính
NEXT\_HOP có thể tiếp tục giữ
địa chỉ của router eBGP bên ngoài.

Điều này có thể khiến các router
nội bộ không truy cập được next-hop
nếu mạng liên kết đến AS bên ngoài
không được quảng bá vào IGP.

Để giải quyết, nhóm sử dụng
lệnh \texttt{next-hop-self}
trên các phiên iBGP.

Ví dụ, khi B1 nhận một tuyến
từ A2 qua eBGP, B1 có thể
quảng bá tuyến đó cho B2
với NEXT\_HOP được thay thành
địa chỉ Loopback0 của B1.

Do các Loopback0 đã được quảng bá
trong OSPF, B2 có thể tìm được
đường đi đến next-hop này.

Tương tự, trong AS A,
A2 và A3 sử dụng
\texttt{next-hop-self}
để các tuyến học từ ISP
có next-hop phù hợp với
khả năng định tuyến nội bộ.

Phương pháp này giúp tách biệt
việc định tuyến nội bộ bằng
RIPv2 hoặc OSPF với quá trình
trao đổi tuyến giữa các AS
thông qua BGP.


% ==========================================
% C. KIEM CHUNG KET QUA
% ==========================================

\subsection{Kiểm chứng kết quả}

\subsubsection{Kiểm tra trạng thái phiên BGP}

Sử dụng lệnh sau trên sáu router
A2, A3, B1, B2, B3 và C1:

\begin{lstlisting}
show ip bgp summary
\end{lstlisting}

Lệnh này cung cấp danh sách
các BGP Neighbor, số AS,
trạng thái phiên và số lượng
prefix được nhận.

Khi phiên BGP đã thiết lập thành công,
cột \texttt{State/PfxRcd} thường
hiển thị số lượng prefix nhận được
thay vì trạng thái
\texttt{Idle}, \texttt{Active}
hoặc \texttt{Connect}.

Các phiên cần kiểm chứng:

\begin{table}[H]
\centering
\caption{Các phiên BGP cần kiểm chứng}
\label{tab:partD-sessions}
\small
\begin{tabular}{cccc}
\toprule
\textbf{Router 1} &
\textbf{Router 2} &
\textbf{Loại} &
\textbf{Yêu cầu} \\
\midrule
A2 & B1 & eBGP & Established \\
A3 & C1 & eBGP & Established \\
B3 & C1 & eBGP & Established \\
A2 & A3 & iBGP & Established \\
B1 & B2 & iBGP & Established \\
B1 & B3 & iBGP & Established \\
B2 & B3 & iBGP & Established \\
\bottomrule
\end{tabular}
\end{table}

% Anh bang chung BGP Summary
% Doi ten anh theo file thuc te tren Overleaf

\partzeroimage
{images/configuration/part_D/A2_summary.png}
{Trạng thái BGP Neighbor trên router A2}
{fig:partD-A2-summary}

\partzeroimage
{images/configuration/part_D/A3_summary.png}
{Trạng thái BGP Neighbor trên router A3}
{fig:partD-A3-summary}

\partzeroimage
{images/configuration/part_D/B1_summary.png}
{Trạng thái BGP Neighbor trên router B1}
{fig:partD-B1-summary}

\partzeroimage
{images/configuration/part_D/B2_summary.png}
{Trạng thái BGP Neighbor trên router B2}
{fig:partD-B2-summary}

\partzeroimage
{images/configuration/part_D/B3_summary.png}
{Trạng thái BGP Neighbor trên router B3}
{fig:partD-B3-summary}

\partzeroimage
{images/configuration/part_D/C1_summary.png}
{Trạng thái BGP Neighbor trên router C1}
{fig:partD-C1-summary}


% ==========================================
% KIEM TRA BGP TABLE
% ==========================================

\subsubsection{Kiểm tra bảng định tuyến BGP}

Sử dụng lệnh:

\begin{lstlisting}
show ip bgp
\end{lstlisting}

trên router A3 và B2 để kiểm tra
các prefix BGP đã được quảng bá
và nhận từ những AS khác.

Các thông tin cần quan sát gồm:

\begin{itemize}

    \item \textbf{Network:}
    Prefix mạng đích.

    \item \textbf{Next Hop:}
    Địa chỉ next-hop của tuyến BGP.

    \item \textbf{Path:}
    Thuộc tính AS\_PATH, biểu diễn
    chuỗi AS mà thông tin tuyến
    đã đi qua.

    \item \textbf{Status:}
    Ký hiệu \texttt{*} cho tuyến hợp lệ
    và \texttt{>} cho tuyến tốt nhất
    được BGP lựa chọn.

\end{itemize}

Các mạng LAN cần xuất hiện
trong quá trình kiểm chứng gồm:

\begin{table}[H]
\centering
\caption{Các mạng LAN được quảng bá vào BGP}
\label{tab:partD-networks}
\small
\begin{tabularx}{\textwidth}{c X X}
\toprule
\textbf{LAN} &
\textbf{Prefix} &
\textbf{Router quảng bá} \\
\midrule
LAN-A & 10.171.126.0/26 & A2, A3 \\
LAN-B & 10.171.126.64/26 & B2 \\
LAN-C & 10.171.126.128/26 & C1 \\
LAN-D & 10.171.126.192/26 & B2, C1 \\
\bottomrule
\end{tabularx}
\end{table}

\partzeroimage
{images/configuration/part_D/A3_bgp.png}
{Bảng BGP trên router A3}
{fig:partD-A3-bgp}

\partzeroimage
{images/configuration/part_D/B2_bgp.png}
{Bảng BGP trên router B2}
{fig:partD-B2-bgp}


% ==========================================
% KIEM TRA KET NOI 4 LAN
% ==========================================

\subsubsection{Kiểm tra kết nối giữa các mạng LAN}

Sau khi các phiên BGP được thiết lập,
thực hiện kiểm tra kết nối giữa
bốn mạng LAN-A, LAN-B, LAN-C và LAN-D.

Việc kiểm tra sử dụng địa chỉ
Loopback1 làm địa chỉ nguồn
để mô phỏng lưu lượng phát sinh
từ những mạng LAN.

Các địa chỉ nguồn và đích
được sử dụng:

\begin{table}[H]
\centering
\caption{Địa chỉ dùng để kiểm tra kết nối LAN}
\label{tab:partD-lan-address}
\small
\begin{tabular}{ccc}
\toprule
\textbf{LAN} &
\textbf{Router} &
\textbf{IP Loopback1} \\
\midrule
LAN-A & A1 & 10.171.126.1 \\
LAN-B & B2 & 10.171.126.65 \\
LAN-C & C1 & 10.171.126.129 \\
LAN-D & D1 & 10.171.126.193 \\
\bottomrule
\end{tabular}
\end{table}

Tổng cộng có sáu cặp LAN
cần được kiểm tra:

\begin{table}[H]
\centering
\caption{Danh sách kiểm tra kết nối giữa các LAN}
\label{tab:partD-lan-ping}
\small
\begin{tabularx}{\textwidth}{c c X c}
\toprule
\textbf{STT} &
\textbf{Nguồn} &
\textbf{Đích} &
\textbf{Kết quả} \\
\midrule
1 & LAN-A & LAN-B & Thành công \\
2 & LAN-A & LAN-C & Thành công \\
3 & LAN-A & LAN-D & Thành công \\
4 & LAN-B & LAN-C & Thành công \\
5 & LAN-B & LAN-D & Thành công \\
6 & LAN-C & LAN-D & Thành công \\
\bottomrule
\end{tabularx}
\end{table}

Ví dụ, để kiểm tra kết nối
từ LAN-A đến LAN-B trên router A1:

\begin{lstlisting}
ping 10.171.126.65 source Loopback1
\end{lstlisting}

Lệnh này gửi gói ICMP đến LAN-B
với địa chỉ nguồn là Loopback1
của router A1.

Các cặp LAN còn lại được kiểm tra
theo phương pháp tương tự.

% Anh ket qua ping 6 cap LAN

\partzeroimage
{images/configuration/part_D/LAN_A_BCD.png}
{Kiểm tra kết nối giữa LAN-A tới LAN-B, LAN-C và LAN-D}
{fig:partD-lan-A-B}


\partzeroimage
{images/configuration/part_D/LAN_B_CD.png}
{Kiểm tra kết nối giữa LAN-B tới LAN-C và LAN-D}
{fig:partD-lan-B-C}

\partzeroimage
{images/configuration/part_D/LAN_C_D.png}
{Kiểm tra kết nối giữa LAN-C và LAN-D}
{fig:partD-lan-C-D}
